% =============================================================================
% introduction.tex — Введение
% =============================================================================

\section*{ВВЕДЕНИЕ}
\addcontentsline{toc}{section}{ВВЕДЕНИЕ}

\textbf{Актуальность темы исследования.}
Ключевая проблема современной быстрой разработки~--- расхождение между кодом и документацией. Темпы внесения изменений в программный продукт настолько высоки, что документирование требований физически не успевает за ними. В результате технические задания становятся неактуальными практически в момент их написания. Причиной тому служат не только хотфиксы, но и реализация устных договорённостей, которые идут вразрез с исходной постановкой.

В результате документация перестаёт быть зеркалом системы: она фрагментарна, неактуальна и зачастую воспринимается разработчиками как формальность. Это приводит к накоплению технического долга, усложняет внедрение в процесс разработки новых специалистов, снижает прозрачность архитектуры сервиса и системы в целом, повышает риск ошибок при доработке или модернизации системы.

При этом проблема признана не только в индустрии, но и в академическом сообществе. Исследования последних лет отмечают систематический разрыв между исходным кодом, архитектурными решениями и документацией~\cite{zhang2022comments,li2023coevolution}. Отдельные работы указывают, что значительная часть проектной документации (до 60--70\%) теряет актуальность спустя всего несколько месяцев после начала разработки, если её поддержка не автоматизирована~\cite{forward2002relevance,storey2021future}.

На этом фоне растёт интерес к использованию больших языковых моделей (LLM) для автоматизированной генерации, обновления и анализа документации по коду. Однако остаётся открытым вопрос: можно ли создать систему, в которой документация не <<догоняет>> код, а рождается вместе с ним и остаётся актуальной автоматически?

\textbf{Цель исследования}~--- разработка и реализация системы автоматической генерации технической документации на основе графовой модели кода и больших языковых моделей, обеспечивающей актуальность и полноту документации при изменениях в кодовой базе.

\textbf{Задачи исследования:}
\begin{enumerate}
  \item Провести анализ существующих подходов к автоматической генерации документации по коду и выявить их ограничения.
  \item Разработать формальную модель представления исходного кода в виде графа знаний, включающего структурные, вызовные и интеграционные связи.
  \item Спроектировать архитектуру системы на основе принципов предметно-ориентированного проектирования (DDD) с модульной организацией.
  \item Реализовать конвейер обработки данных: от клонирования репозитория до генерации документации с помощью LLM.
  \item Разработать подсистему RAG (Retrieval-Augmented Generation) с графовым расширением контекста для интерактивного взаимодействия с системой.
  \item Реализовать визуализацию графа кода и кросс-приложенческих интеграций.
  \item Провести тестирование и оценку качества генерируемой документации.
\end{enumerate}

\textbf{Объект исследования}~--- процесс автоматической генерации и поддержания актуальности технической документации программных систем.

\textbf{Предмет исследования}~--- методы и алгоритмы построения графовых моделей кода, семантического чанкования, векторного поиска и многоступенчатой генерации документации с помощью LLM.

\textbf{Научная новизна:}
\begin{enumerate}
  \item Предложена унифицированная графовая модель представления кода, включающая 24~типа узлов и 24~типа связей, покрывающая структурные, вызовные и интеграционные отношения между компонентами программной системы.
  \item Разработан метод многоступенчатой генерации документации (Coder~$\to$~Talker~$\to$~Digest), адаптирующий описания под различные целевые аудитории~--- разработчиков и бизнес-аналитиков.
  \item Предложен подход к автоматическому обнаружению кросс-приложенческих интеграций (HTTP, Kafka, Camel) на основе статического анализа кода и байткода.
  \item Реализован метод RAG с графовым расширением контекста, комбинирующий обход графа зависимостей и векторный поиск для ответов на вопросы о коде.
\end{enumerate}

\textbf{Практическая значимость} работы определяется реализацией полнофункциональной системы, пригодной к внедрению в промышленную разработку. Система позволяет автоматизировать документирование корпоративных приложений на Kotlin/Java, визуализировать архитектуру и интеграции, а также предоставляет интерактивный интерфейс для получения ответов о коде.

\textbf{Методы исследования:} статический анализ исходного кода (PSI, ASM), графовое моделирование, векторные представления (embeddings), генеративные языковые модели, предметно-ориентированное проектирование (DDD), тестирование программного обеспечения.

\textbf{Структура работы.} Работа состоит из введения, четырёх глав, заключения, списка использованных источников и приложений.
